\documentclass{memoir}

\usepackage[spanish,mexico,es-noindentfirst]{babel}
\usepackage{mathtools}
\usepackage{amssymb}
\usepackage{mathrsfs}
\usepackage{mycats}

\begin{document}
\tableofcontents*

\chapter*{Introducción}
Creo que ya hay algo de trabajo en la introducción.

\chapter{Preliminares}
\section{Conjuntos abstractos}
En esta sección debes escribir los axiomas elementales de la categoría de
conjuntos. También estaría bien que esté la discusión sobre su relación con ZFC.

\section{Comónadas y coálgebras}
Dar aquí la definición de comónada. Después hay que observar que toda adjunción
induce una comónada. Luego definir la categoría de coálgebras y notar que hay un
funtor que olvida. Después de esto dar la coálgebra libre y notar que esto
induce un funtor que es adjunto izquierdo del funtor que olvida.

Con lo anterior se tiene una situación como la siguiente
\begin{equation*}
  \begin{tikzcd}
    \catB\ar[shift right]{d}[swap]{L} & \catA_C\ar[shift right]{d}[swap]{U}\\
    \catA\ar[shift right]{u}[swap]{R}\ar[equals]{r} 
     & \catA\mathrlap{.}\ar[shift right]{u}[swap]{F}
  \end{tikzcd}
\end{equation*}
La unidad y la counidad de \(L\dashv R\) serán denotadas con
\(\alpha\) y \(\beta\), respectivamente. En la situación del diagrama es posible
construir un funtor de comparación \(K\colon\catB\to\catA_C\) que conmute con
los adjuntos izquierdos y por tanto con los derechos. De hecho este funtor es
único, pero no sé si vale la pena mostrar la unicidad. Para que conmute con los
derechos debe cumplir que \(KB=(LB,b)\), donde \(b\colon LB\to LRLB\) es una
estructura de coálgebra. El valor que podría tener hacer la unicidad de \(K\) es
que se muestra que la estructura debe ser \(L\alpha_B\colon LB\to LRLB\). La
otra es que se proponga tal cual y se muestre que hace conmutar al diagrama de
los izquierdos. Este funtor de comparación es el funtor que aparece en la
demostración del teorema~2.27. Por lo tanto, tal vez valga la pena comemtar que
no es necesario demostrar el dual del teorema de Beck, ya que en el caso que te
interesa se puede ver fácilmente que sí es una equivalencia.

Finalmente, terminar con la demostración de \(\toposE_C\) es un topos.

\section{Objetos decidibles}
Dar aquí las propiedades de decidibles que no usan a la categoría
\(\toposM\). Esto es que los subobjetos de decidibles son decidibles y que el
terminal y el objeto de números naturales lo son. Con esto las cosas de
naturales también estarán aquí.

Si puedes escribe algunas propiedades más de los artículos de decidibles que
comentamos al principio.

\chapter{El topos de espacios}
\section{Motivación y definición o algo así}
Escribir aquí los axiomas que harán que consideremos que un topos sea
considerado una categoría de espacios. Decir que estos axiomas están motivados
por lo que cumple el topos de la geometría diferencial sintética.

\section{El subtopos de espacios discretos o decidibles}
Obtener al subtopos \(\toposM\) y los dos funtores adjuntos \(\Gamma\) y
\(\Delta\). Con esto desde la proposición~2.24 en adelante estará en este
capítulo.

\section{El concepto de espacio como noción fundamental}
En esta parte mostrar que \(\toposM\) satisface los axiomas elementales de la
categoría de conjuntos. Concluir que la noción de conjunto es derivada de la de
espacio.



\end{document}